% Revize 2026-10-28: obsah, matematika a návaznost; podrobnosti a časy
% viz planovani/REVIZE_PRIPRAV_1_ROCNIK.md.
% Aktuální porovnání se skutečným obsahem místních PDF Realisticky.cz:
% externi_zdroje/realisticky/01_Zakladni_poznatky/06_Mocniny/010610_Mocniny_shrnuti/010610_Lekce.pdf
% externi_zdroje/realisticky/01_Zakladni_poznatky/02_Ciselne_obory/010209_Usmernovani_zlomku/010209_Lekce.pdf
% Revize 2026-10-27: cílené doplnění; nové úlohy a řešení jsou vlastní tvorba.
% Metodický vzor: 4_rocnik/pripravy/Uvod_do_kombinatoriky.pdf (samostatný dokument).
% Hodina 39, 1. ročník SPŠ. Rozsah: Rozpis_vyuky_MAT_1_rocnik.tex, řádek 39.

% Plán: 1_rocnik/tematicke_plany/SPS.pdf; společná projektová šablona.

% Učebnice: ucebnice/01_Zakladni_poznatky/Mocniny_a_odmocniny/Mocniny a Odmocniny.tex

% Původ úloh: Vlastní zadání a vlastní formulace výkladu; metodická inspirace zdroji.

% Metodický zdroj: externi_zdroje/realisticky/01_Zakladni_poznatky/06_Mocniny/010610_Mocniny_shrnuti/010610_Lekce.pdf

% URL: http://www.realisticky.cz/ucebnice/01%20Matematika%20S%C5%A0/01%20Z%C3%A1kladn%C3%AD%20poznatky/06%20Mocniny/10%20Mocniny%20-%20shrnut%C3%AD.pdf

% Metodický zdroj: externi_zdroje/realisticky/01_Zakladni_poznatky/02_Ciselne_obory/010209_Usmernovani_zlomku/010209_Lekce.pdf

% URL: http://www.realisticky.cz/ucebnice/01%20Matematika%20S%C5%A0/01%20Z%C3%A1kladn%C3%AD%20poznatky/02%20%C4%8C%C3%ADseln%C3%A9%20obory/09%20Usm%C4%9Br%C5%88ov%C3%A1n%C3%AD%20zlomk%C5%AF.pdf

\documentclass[11pt,a4paper]{article}

% Spolecny styl pro ucitelske pripravy MAT (Work/MAT)
% Verze 1.1, 2026-10-23
\NeedsTeXFormat{LaTeX2e}
\ProvidesPackage{priprava}[2026/10/23 v1.1 Sablona ucitelskych priprav MAT]

% Zakladni typografie a cestina
\RequirePackage[T1]{fontenc}
\RequirePackage[utf8]{inputenc}
\RequirePackage[czech]{babel}
\RequirePackage{lmodern}
\RequirePackage{microtype}
\RequirePackage{amsmath,amssymb}
\RequirePackage{xcolor}
\RequirePackage{enumitem}
\RequirePackage[most]{tcolorbox}
\RequirePackage{tabularx,booktabs,array}
\RequirePackage{fancyhdr}
\RequirePackage{needspace}
\RequirePackage{geometry}
\RequirePackage[hidelinks]{hyperref}

\geometry{a4paper,top=17mm,bottom=18mm,left=23mm,right=23mm}
\definecolor{notegray}{gray}{0.39}
\setlength{\parindent}{0pt}
\setlength{\parskip}{2.5pt plus 1pt minus 0.5pt}
\setlist[itemize]{topsep=2pt,itemsep=1pt,parsep=0pt,leftmargin=1.45em}
\setlist[enumerate]{topsep=2pt,itemsep=2pt,parsep=0pt,leftmargin=1.65em}
\setcounter{secnumdepth}{0}
\renewcommand{\arraystretch}{1.12}

% Zahlavi: vlevo rocnik | globalni tema, vpravo cislo: nazev hodiny
% Zapati zustava zcela prazdne.
\pagestyle{fancy}
\fancyhf{}
\renewcommand{\headrulewidth}{0.35pt}
\renewcommand{\footrulewidth}{0pt}
\setlength{\headheight}{14pt}
\newcommand{\nastavPripravu}[4]{%
  \fancyhead[L]{\small #1\enspace\textbar\enspace #2}%
  \fancyhead[R]{\small #3: #4}%
}

% Minimalisticke cernobile boxy. Kazda definice zvlast.
\tcbset{commonlesson/.style={enhanced,breakable=false,colback=white,colframe=black,
  colbacktitle=white,coltitle=black,boxrule=0.45pt,
  arc=0.8mm,outer arc=0.8mm,left=3mm,right=3mm,top=1.6mm,bottom=1.6mm,
  toptitle=1mm,bottomtitle=0.7mm,fonttitle=\bfseries,
  attach boxed title to top left={xshift=2.5mm,yshift=-1.15mm},
  boxed title style={colback=white,colframe=white,boxrule=0pt,left=1mm,right=1mm}}}
\newtcolorbox{zadaniBox}[1][]{commonlesson,title={Zadání},#1}
\newtcolorbox{prikladBox}[1][]{commonlesson,title={Příklad},#1}
\newtcolorbox{definiceBox}[1][]{commonlesson,title={Definice},#1}

% Spolecny box pro zadani a jeho reseni, bez mezery mezi castmi.
% Pouziti:
%   \begin{prikladReseniBox}
%     Zadani ulohy.
%     \tcblower
%     \begin{align*} ... \end{align*}
%   \end{prikladReseniBox}
% Spodni cast ma jemny sedy podklad a delici linku; lze v ni sazet
% vice radku matematiky. Obe casti maji jen jeden spolecny ramecek.
\tcbset{solutionlesson/.style={
  commonlesson,
  enhanced jigsaw,
  breakable,
  colbacklower=black!3,
  segmentation style={solid,draw=black!25,line width=0.35pt},
  fontlower=\small\color{notegray},
  before lower={\noindent\textbf{Řešení:}\enspace},
  leftlower=3mm,rightlower=3mm
}}
\newtcolorbox{prikladReseniBox}[1][]{solutionlesson,title={Příklad},#1}
\newtcolorbox{zadaniReseniBox}[1][]{solutionlesson,title={Zadání},#1}

% Pomucky pro autora pripravy: sedive poznamky; \reseni pro starsi pripravy
\newcommand{\poznamka}[1]{\par\begingroup\small\color{notegray}\itshape
  \textnormal{\textbf{Poznámka pro učitele.}}\enspace #1\par\endgroup}
\newcommand{\reseni}[1]{\par\begingroup\small\color{notegray}
  \textbf{Očekávané řešení:}\enspace #1\par\endgroup}
\newcommand{\blok}[1]{\needspace{5\baselineskip}\subsection*{#1}}
\newcommand{\vyrok}[1]{\uv{#1}}


\nastavPripravu{1. ročník}{Odmocniny}{39}{Opakování mocnin a odmocnin}

\begin{document}

\blok{Výklad a souvislosti}

Hodina propojuje již známá pravidla. U každého výrazu nejprve pojmenuj vnější operaci, určuj podmínky a vol úpravu, která zápis zjednoduší. Výsledky ověřuj zpětným dosazením nebo umocněním; jediná číselná zkouška však nedokazuje obecnou rovnost.

\begin{definiceBox}[title={Připomenutí definice},breakable=false]
\textbf{Racionální mocnina} kladného $a$ je pro celá $m$ a přirozená $n\ge2$ dána vztahem $a^{m/n}=(\sqrt[n]a)^m$.
\end{definiceBox}


Připomeň rozdíl $(-a)^n$ a $-a^n$, nenulové základy u záporných mocnin, nezáporné výsledky druhých odmocnin a zachování podmínek při krácení. Zásobu rozděl mezi společný rozbor a samostatnou práci podle skutečných chyb; nejde o nový tematický blok.

\blok{Řešené příklady}

\begin{prikladReseniBox}[title={Příklad 1},breakable=false]
Pro $n\in\mathbb Z$ uprav $\frac{3^{2n+1}\cdot9^{n-1}}{27^{n-2}}$.
\tcblower
Převodem na základ $3$ získáme $3^{2n+1+2n-2-3n+6}=3^{n+5}$. Jmenovatel není pro žádné celé $n$ nulový.
\end{prikladReseniBox}


\begin{prikladReseniBox}[title={Příklad 2},breakable=false]
Uprav $\frac{\sqrt{12}+\sqrt{27}}{\sqrt3}$.
\tcblower
Po částečném odmocnění dostaneme $\frac{2\sqrt3+3\sqrt3}{\sqrt3}=5$. Jmenovatel je kladný, krácení je přípustné.
\end{prikladReseniBox}


\blok{Samostatné procvičení}

\begin{zadaniReseniBox}[title={Procvičení},breakable=false]
\begin{enumerate}[label=\alph*)]
\item Uprav $(2a^{-2}b^3)^2$ a uveď podmínky.
\item Vypočítej $16^{-3/4}$ a uprav $\sqrt{200}$.
\item Usměrni $\frac1{\sqrt5+2}$.
\end{enumerate}
\tcblower
\begin{enumerate}[label=\alph*)]
\item $\frac{4b^6}{a^4}$ při $a\ne0$; $b$ může být i nula.
\item $1/8$ a $10\sqrt2$.
\item $\sqrt5-2$, protože $(\sqrt5+2)(\sqrt5-2)=1$.
\end{enumerate}
\end{zadaniReseniBox}


\poznamka{U procvičení a) výslovně porovnej podmínku $a\ne0$ s nepotřebným zákazem $b=0$. V záporném exponentu se nenachází $b$. Při pomalejším tempu vynech rezervu a zopakuj původní chybné kroky žáků.}


\blok{Navazující procvičování ve třídě}

\begin{zadaniReseniBox}[title={Rezerva: oprav chybné řešení},breakable=false]
\begin{enumerate}[label=\alph*)]
\item Žák píše $\sqrt{9a^2}=3a$ pro všechna reálná $a$. Oprav výsledek a vyvrať jeho tvrzení.
\item Žák píše $\frac{a^{-2}}{a^{-5}}=a^{-7}$. Oprav výpočet i podmínky.
\item Uprav $\frac{\sqrt{48}+\sqrt{75}-\sqrt{27}}{\sqrt3}$.
\end{enumerate}
\tcblower
\begin{enumerate}[label=\alph*)]
\item Správně $3|a|$. Pro $a=-1$ je levá strana $3$, ale navržený výsledek $-3$.
\item $a^{-2-(-5)}=a^3$ pro $a\ne0$.
\item $\frac{4\sqrt3+5\sqrt3-3\sqrt3}{\sqrt3}=6$.
\end{enumerate}
\end{zadaniReseniBox}

\blok{Volitelná rezerva}

\begin{prikladReseniBox}[title={Rozšíření procvičení},breakable=false]
Pro reálná $x$ uprav $\frac{\sqrt{50x^2}}{5x}+\frac1{\sqrt{x-1}}$ a urči obor.
\tcblower
První jmenovatel vyžaduje $x\ne0$, druhý $x-1>0$. Společně $x>1$, tedy $|x|=x$. Dostaneme $\sqrt2+\frac{\sqrt{x-1}}{x-1}$ pro $x>1$.
\end{prikladReseniBox}



\poznamka{Základ tvoří původní opakovací skupiny. Při svižném tempu následuje nová diagnostická rezerva: žák označí první chybný krok a zdůvodní opravu. Domácí čtveřice je procvičování, nikoli návrh testu. Není třeba znovu opisovat všechny definice; čas věnuj vysvětlení volby pravidla.}

\blok{Domácí procvičení}

\begin{zadaniReseniBox}[title={Domácí opakování před průběžným ověřením},breakable=false]
\begin{enumerate}[label=\alph*)]
\item Uprav $\frac{(2a^{-1}b^2)^3}{4a^{-2}b}$ a stanov podmínky.
\item Vypočítej $27^{-2/3}$ a $\sqrt{108}$.
\item Usměrni $\frac2{\sqrt6-2}$.
\item Pro $n\in\mathbb Z$ uprav $\frac{2^{n+2}\cdot4^n}{8^{n-1}}$.
\end{enumerate}
\tcblower
\begin{enumerate}[label=\alph*)]
\item $2b^5/a$, $a,b\ne0$.
\item $1/9$ a $6\sqrt3$.
\item $\sqrt6+2$, protože sdružený součin je $2$.
\item $2^{n+2+2n-3n+3}=2^5=32$.
\end{enumerate}
\end{zadaniReseniBox}

\end{document}
